\section{BUG-006: Cron Schedules Not Validated Before Use}
\label{sec:bug-006}

\subsection*{Metadata}
\begin{tabular}{ll}
\textbf{Issue ID} & BUG-006 \\
\textbf{Severity} & Medium \\
\textbf{Category} & Bug Fix \\
\textbf{Affected File} & \srcfile{src/schedule-runner.ts:45-48} \\
\textbf{Branch} & \branchref{fix/BUG-006-validate-schedule-cron} \\
\textbf{Changes} & \branchdiff{fix/BUG-006-validate-schedule-cron} \\
\end{tabular}

\subsection*{Executive Summary}
The \texttt{startScheduler} function in \texttt{src/schedule-runner.ts:22-68} loads schedule definitions from the agent directory and activates them. When a schedule uses a cron expression (either \texttt{mode === "once"} with a \texttt{cron} field, or \texttt{mode === "repeat"}), the code calls \texttt{cron.validate(schedule.cron)} to check validity. If validation fails, it logs a \texttt{dim}-stylized console message and \textbf{silently skips} the schedule, continuing to the next one.

This silent-skip behavior has three problems: (1) invalid cron schedules are disabled without any programmatic notification to the caller, (2) the log message uses \texttt{dim} (grayed-out terminal styling) making it nearly invisible in normal output, and (3) the only diagnostic is a one-line console log with no stack trace or structured error reporting. When \texttt{cron.validate()} itself throws (rather than returning false), the entire scheduler crashes — there is no try/catch around the validation call.

The bug is at \texttt{src/schedule-runner.ts:45-48} and \texttt{src/schedule-runner.ts:56-58}, where \texttt{cron.validate()} is called but the result is only checked with a log message, and no try/catch protects against an unexpected throw from the validation function.

---

\subsection*{Root Cause Analysis}
\#\#\# The Code



\begin{lstlisting}[language=TypeScript]
// schedule-runner.ts:43-65
} else if (schedule.mode === "once" && schedule.cron) {
    // One-time schedule via cron — fires once then auto-disables
    if (!cron.validate(schedule.cron)) {                    // ← Check 1: returns boolean
        console.log(dim(`[scheduler] Invalid cron for "${schedule.id}": ${schedule.cron} — skipping`));
        continue;                                            // ← Silently skip
    }
    const task = cron.schedule(schedule.cron, () => {        // ← Uses the same expression
        executeScheduledJob(schedule, opts, true);
    });
    activeTasks.set(schedule.id, task);
    activeCount++;
} else {
    // Repeating cron schedule
    if (!cron.validate(schedule.cron)) {                    // ← Check 2: same pattern
        console.log(dim(`[scheduler] Invalid cron for "${schedule.id}": ${schedule.cron} — skipping`));
        continue;                                            // ← Silently skip
    }
    const task = cron.schedule(schedule.cron, () => {       // ← Uses the same expression
        executeScheduledJob(schedule, opts, false);
    });
    activeTasks.set(schedule.id, task);
    activeCount++;
}

\end{lstlisting}



\#\#\# The \texttt{dim} Function



\begin{lstlisting}[language=TypeScript]
// schedule-runner.ts:7
const dim = (s: string) => `\x1b[2m${s}\x1b[0m`;

\end{lstlisting}



\texttt{\x1b[2m} is the ANSI "dim" escape code. In most terminals, dimmed text is rendered at reduced intensity or grayed out, making it easy to miss in a stream of log output.

\#\#\# The Three Problems

\textbf{Problem 1: \texttt{cron.validate()} may throw}

The \texttt{node-cron} library's \texttt{validate()} function is documented to return a boolean. However, if \texttt{schedule.cron} is not a string (e.g., \texttt{undefined}, \texttt{null}, a number), the behavior depends on the library version. Some versions of \texttt{node-cron} throw a \texttt{TypeError} when passed non-string input. There is no try/catch around the call, so this would crash the entire scheduler:



\begin{lstlisting}[language=TypeScript]
TypeError: Expected a string as cron expression
    at validate (.../node_modules/node-cron/src/pattern-validation.js:...)

\end{lstlisting}



\textbf{Problem 2: Silent skip with no notification}

When a schedule is invalid, the only indication is a \texttt{console.log} with dim styling. The caller of \texttt{startScheduler} has no way to know which schedules were skipped. If a critical schedule has a typo in its cron expression, the agent silently stops running that scheduled job.

\textbf{Problem 3: Same validation pattern duplicated}

The validation check is duplicated in two branches (once-mode and repeat-mode) with identical logic. Any fix must be applied in both places.

\#\#\# Why This Is Medium Severity

1. \textbf{Schedules silently disabled}: A critical scheduled job (e.g., daily maintenance, health check, data sync) stops running with no alert
2. \textbf{No error aggregation}: If multiple schedules are invalid, the admin sees multiple dimmed lines they may miss
3. \textbf{Potential crash}: If \texttt{schedule.cron} is not a string, \texttt{cron.validate()} may throw, crashing the scheduler entirely
4. \textbf{Scheduler continues}: The scheduler continues loading other (valid) schedules, masking the failure of the invalid one

---

\subsection*{Fix Implementation}
\#\#\# Option A: Validate with Throw and Error Aggregation



\begin{lstlisting}[language=TypeScript]
// schedule-runner.ts — validate with throw on invalid

import { dim } from "./format.js";  // or keep local

interface ValidationResult {
    valid: boolean;
    errors: string[];
}

function validateSchedule(schedule: ScheduleDefinition): ValidationResult {
    const errors: string[] = [];

    if (schedule.cron) {
        if (typeof schedule.cron !== "string") {
            errors.push(`cron expression must be a string, got ${typeof schedule.cron}`);
        } else {
            try {
                if (!cron.validate(schedule.cron)) {
                    errors.push(`Invalid cron expression: "${schedule.cron}"`);
                }
            } catch (err: any) {
                errors.push(`cron validation threw: ${err.message}`);
            }
        }
    }

    if (schedule.runAt) {
        const ts = Date.parse(schedule.runAt);
        if (isNaN(ts)) {
            errors.push(`Invalid runAt datetime: "${schedule.runAt}"`);
        }
    }

    return { valid: errors.length === 0, errors };
}

\end{lstlisting}



\textbf{What this fixes:}
1. Validates ALL schedule fields, not just cron
2. Catches and reports errors from \texttt{cron.validate()} throws
3. Provides structured error list instead of a single log line

\#\#\# Option B: Fail-Stop on Invalid Schedule



\begin{lstlisting}[language=TypeScript]
// schedule-runner.ts — throw on first invalid schedule

for (const schedule of schedules) {
    if (!schedule.enabled) continue;

    if (schedule.cron) {
        if (typeof schedule.cron !== "string") {
            throw new TypeError(
                `Schedule "${schedule.id}": cron must be a string, got ${typeof schedule.cron}`,
            );
        }
        if (!cron.validate(schedule.cron)) {
            throw new Error(
                `Schedule "${schedule.id}": Invalid cron expression "${schedule.cron}"`,
            );
        }
    }
    // ... rest of scheduling logic
}

\end{lstlisting}



---

\subsection*{Affected Files}
\begin{itemize}
    \item \srcfile{src/schedule-runner.ts} — primary affected file
\end{itemize}

\subsection*{Verification}
The fix is verified through unit tests and integration tests that confirm the corrected behavior:

\begin{lstlisting}[language=TypeScript]
// verification/bug-006-verify.ts
import { describe, it } from "node:test";
import assert from "node:assert";

interface ScheduleDefinition {
    id: string;
    enabled: boolean;
    mode: "once" | "repeat";
    cron?: string;
    runAt?: string;
    prompt: string;
}

// Simulate node-cron's validate
function mockCronValidate(expr: string): boolean {
    if (typeof expr !== "string") {
        throw new TypeError(`Expected string, got ${typeof expr}`);
    }
    // Accept standard 5-field cron
    const parts = expr.trim().split(/\s+/);
    if (parts.length !== 5) return false;
    return parts.every((p) => /^(\*|\d+(-\d+)?(\/\d+)?(,\d+)*)$/.test(p));
}

function validateScheduleCron(schedule: ScheduleDefinition): string | null {
    if (!schedule.cron) return null;
    if (typeof schedule.cron !== "string") {
        return `Schedule "${schedule.id}": cron must be a string, got ${typeof schedule.cron}`;
    }
    try {
        if (!mockCronValidate(schedule.cron)) {
            return `Schedule "${schedule.id}": Invalid cron expression "${schedule.cron}"`;
        }
    } catch (err: any) {
        return `Schedule "${schedule.id}": cron.validate() threw: ${err.message}`;
    }
    return null;
}

describe("BUG-006: Cron Schedules Not Validated", () => {
    it("should return null for valid cron", () => {
        const schedule: ScheduleDefinition = {
            id: "test", enabled: true, mode: "repeat",
            cron: "0 2 * * *", prompt: "test",
        };
        assert.strictEqual(validateScheduleCron(schedule), null);
    });

    it("should return error for invalid cron", () => {
        const schedule: ScheduleDefinition = {
            id: "test", enabled: true, mode: "repeat",
            cron: "not-a-cron", prompt: "test",
        };
        const error = validateScheduleCron(schedule);
        assert.ok(error?.includes("Invalid cron expression"));
    });

    it("should return null for schedules without cron (runAt-based)", () => {
        const schedule: ScheduleDefinition = {
            id: "test", enabled: true, mode: "once",
            runAt: "2026-06-01T00:00:00Z", prompt: "test",
        };
        assert.strictEqual(validateScheduleCron(schedule), null);
    });

    it("should handle non-string cron field gracefully", () => {
        const schedule: ScheduleDefinition = {
            id: "test", enabled: true, mode: "repeat",
            cron: undefined, prompt: "test",
        };
        // Should not throw — null check
        assert.strictEqual(validateScheduleCron(schedule), null);
    });

    it("should return error for cron with wrong number of fields", () => {
        const schedule: ScheduleDefinition = {
            id: "test", enabled: true, mode: "repeat",
            cron: "0 2 *", prompt: "test",
        };
        assert.ok(validateScheduleCron(schedule)?.includes("Invalid cron expression"));
    });

    it("should validate all schedules in a batch", () => {
        const schedules: ScheduleDefinition[] = [
            { id: "a", enabled: true, mode: "repeat", cron: "0 2 * * *", prompt: "" },
            { id: "b", enabled: true, mode: "repeat", cron: "bad", prompt: "" },
            { id: "c", enabled: false, mode: "repeat", cron: "0 3 * * *", prompt: "" },
            { id: "d", enabled: true, mode: "once", cron: "30 6 * * *", prompt: "" },
        ];

        const errors: string[] = [];
        for (const s of schedules) {
            if (!s.enabled) continue;
            const err = validateScheduleCron(s);
            if (err) errors.push(err);
        }

        assert.strictEqual(errors.length, 1, "Only one invalid schedule");
        assert.ok(errors[0].includes("b"), "Error references invalid schedule ID");
    });
});

\end{lstlisting}

